\scanpage{86}
% loc 1–7
le système $(b,c,b',c',t)$ ``pour presque tout \deleted{choi} $t$'' \inserted{et loc fpqc} par deux éléments \inserted{$u,v$} \unsure{de dans|dedans|donnant} l'alg de Matrices \inserted{$A=M_2(k)$}. En effet, \deleted{partant} \deleted{de \unsure{une|$u,v$}}\inserted{quitte} à se localiser pour que $A_1,A_2$ soient \unsure{scindies|scindées} $(\simeq k\times k)$, on peut les réaliser comme plongées dans $A$, i.e. par $u\in A$, $v\in A$.
% loc 8–10
Bien sûr, on ne \deleted{pas de \unsure{en|ce}} peut préjuger de la valeur de $\operatorname{Tr}uv$. Mais \deleted{considérons, pour}\inserted{on a}

% loc 11–14, lemma
\uline{Lemme} Soient $u,v\in A\ (\simeq M_2(k))$ qui \deleted{soient} engendrent $k[u],k[v]$ étales rang 2 sur base $\operatorname{Spec}k$. Considérons pour $\unsure{\theta|\Theta}$ variable dans $A^*$
% loc 15
\[ \operatorname{Tr}(u\theta(v)) \qquad \theta(v)\overset{\text{dfn}}{=}\theta v\theta^{-1} \]
% loc 16
donc on trouve morphisme de schémas
% loc 17
\[ \deleted{\unsure{M_2|M}}\ \check{\mathbb V}(A)\longrightarrow\underbrace{\mathbb E^1_k}_{\text{espace affine}=\operatorname{Spec}k[T]} \]
% loc 17–18, oblique note to left
% slanted
\margin{\uline{NB} $\theta\unsure{\longrightarrow|\longmapsto}\theta(v)$ est épi sur les conjugués à $v$,}
% loc 18
Celui-ci est \uline{dominant} --- 

% loc 19–23
En fait, on trouve une fonction rationnelle en $b,c,b',c',t$, ayant comme dénominateur $t^2-4cc'$ $\bigl(=\deleted{\det}\operatorname{Tr}uv^2-4\det(uv)=\delta(uv)$, discriminant de $uv\bigr)$
% archive sheet number: 85
